\begin{recipebox}{Beef, Red Lentil \& Vegetable Cottage Pie Bowls}{4}{20 min prep + \textasciitilde35 min}

All the good parts of cottage pie without assembling and baking a pie. The
filling and potatoes cook together: potatoes steam in a basket above the mince.
Roughly 40--43\,g protein per serve.

\begin{ingredients}
  \item[] \textbf{Filling}
  \item 500\,g lean beef mince
  \item \nicefrac{3}{4} cup (140--150\,g) dry split red lentils, rinsed
  \item 2 carrots, finely diced
  \item 1 onion, diced
  \item 3 garlic cloves, crushed
  \item 1\nicefrac{1}{2} cups beef stock
  \item 2 tbsp tomato paste
  \item 1 tbsp Worcestershire sauce
  \item 1 tsp dried thyme
  \item 1 tsp smoked paprika
  \item 300\,g frozen peas
  \item[] \textbf{Mash}
  \item 800\,g potatoes, peeled and cut into chunks
  \item \nicefrac{1}{2} cup low-fat milk
  \item Salt and pepper
\end{ingredients}

\begin{method}
  \item \textbf{Soften the veg.} Select Saut\'e and cook the onion and carrots in a little oil for 3--4 minutes.
  \item \textbf{Brown the mince.} Add the mince and cook until browned. Add the garlic, thyme and paprika for the last 30 seconds.
  \item \textbf{Deglaze.} Pour in the stock and scrape the base completely clean. Stir in the Worcestershire sauce.
  \item \textbf{Layer --- don't stir.} Scatter the rinsed lentils evenly over the top so they sit in the liquid rather than settling on the base, then dot the tomato paste over them. \textbf{Do not stir.}
  \item \textbf{Potatoes above.} Set a trivet or steamer basket over the filling and put the potatoes in it.
  \item \textbf{Pressure cook.} Lid on, valve sealed. Cook on High for \textbf{7 minutes}, then a \textbf{10-minute natural release}.
  \item \textbf{Mash.} Lift out the potatoes and mash with the milk, salt and pepper.
  \item \textbf{Finish the filling.} Stir the filling well (the lentils will have softened into the gravy), then stir in the frozen peas. Use Saut\'e for a few minutes if it needs thickening.
  \item Divide the filling between four bowls and spoon the mash over the top.
\end{method}

\begin{notelist}
  \item \textbf{Why the lentils go on top:} red lentils and tomato paste stirred into the base are a classic burn-warning trigger. Layering them over the deglazed mince keeps the base thin while they cook.
  \item Use \textbf{split red lentils} --- canned brown lentils won't thicken the gravy. In a pinch, stir canned lentils in after pressure cooking with the peas.
  \item \textbf{Advance prep:} defrost the mince in the fridge the night before. Peel the potatoes shortly before cooking (or keep them submerged in water in the fridge).
\end{notelist}

\end{recipebox}
